\chapter*{2007 年之后的进展}
\addcontentsline{toc}{chapter}{2007 年之后的进展}
\markboth{2007 年之后的进展}{2007 年之后的进展}

\begin{center}
\large 本书成书于 2007 年；以下为截至 2026 年的补充
\end{center}

\section*{阅读说明}

本章不属于原书，也不改动原书任何章节的编号。写作目的是回答一个读完全书自然会问的问题：书末第 10 章把局部几何 Langlands 对应陈述为猜想，快二十年过去，这个猜想证到了什么程度？哪些还是猜想？哪些方向是本书完全没有预料到的？

本章按“已证 / 部分结果 / 猜测”三级标注。所有条目的来源都给出 arXiv 编号或期刊 DOI，读者可自行核对；凡本轮未能核实的，明确写“未核实”，不把预印本说成已发表，也不把物理论文说成数学证明。

\section*{1. 整体几何 Langlands：2024 年的五篇证明}

\subsection*{1.1 证明了什么}

2024 年 5 月至 9 月，Dennis Gaitsgory、Sam Raskin 与另外七位合作者（Arinkin、Beraldo、Campbell、Chen、Faergeman、Lin、Rozenblyum）以五篇预印本给出了几何 Langlands 猜想在下列情形下的证明：
\begin{itemize}
\item 底曲线是特征 $0$ 的光滑射影曲线；
\item 群是约化群；
\item 局部系统处处无分歧（unramified）。
\end{itemize}
五篇的 arXiv 编号依次为 \texttt{2405.03599}（I，构造函子）、\texttt{2405.03648}（II，Kac--Moody 局部化与 FLE）、\texttt{2409.07051}（III，与抛物诱导的相容性）、\texttt{2409.08670}（IV，ambidexterity）、\texttt{2409.09856}（V，multiplicity one）。合计八百余页。截至 2026 年 11 月，五篇均未在期刊发表，也未完成同行评议；引用时应写“预印本”。

证明的现代形式是 Arinkin--Gaitsgory 在 2015 年给出的范畴等价
\[
D\text{-}\mathrm{mod}(\mathrm{Bun}_G)\;\simeq\;\mathrm{IndCoh}_{\mathcal N}(\mathrm{LS}_{{}^L G}),
\]
其中 $\mathrm{IndCoh}_{\mathcal N}$ 表示在幂零奇异支撑条件下取的 ind-coherent 层。2024 年的五篇证明的正是这个等价。第 I 篇构造自守侧到谱侧的函子；第 V 篇的 multiplicity one 定理说明不可约局部系统对应的 Hecke 本征层在相差一个向量空间的意义下唯一，由此收束反方向。

\subsection*{1.2 没有证明什么}

这一点比“证明了什么”更重要，因为本书读者最容易把 2024 年的结果误当成“整个 Langlands 纲领完成了”。
\begin{enumerate}
\item \textbf{分歧情形没有证。} 五篇只处理处处无分歧的局部系统。分歧情形（本书第 10 章 §10.4 的温和分歧猜想）不在其中。IHES 的官方介绍与 \emph{Nature} 2025 年的报道都明确用了“unramified”这一限定。
\item \textbf{正特征没有证。} 主证明在特征 $0$；向特征 $p$ 的转移是另一篇预印本（Gaitsgory--Raskin，\texttt{arXiv:2508.02237}，2025），不属于这五篇。
\item \textbf{局部几何 Langlands 没有证。} 本书第 10 章的局部猜想（Frenkel--Gaitsgory）不在五篇结论之内。第 II 篇把临界水平 Kac--Moody 局部化当作\emph{工具}，结论仍是整体对应。
\item \textbf{同行评议没有完成。} 截至 2026 年 11 月，五篇的 arXiv 页均无 Journal-ref，Crossref 也无期刊记录。
\end{enumerate}

\subsection*{1.3 与本书的关系}

2024 年的证明走的是纯数学路线：Kac--Moody 局部化、FLE（fundamental local equivalence）、ambidexterity、multiplicity one。它与 Kapustin--Witten 的物理图像互相印证，但不依赖后者。对本书读者而言，最直接的关联是：第 II 篇的局部化工具正是建立在本条第 2--8 章所发展的临界水平表示论之上；换句话说，本书不只是“历史背景”，它的技术被用进了 2024 年的证明。

\section*{2. 局部几何 Langlands：本书第 10 章的猜想现状}

本书 §10.3 的无分歧情形（$G[[t]]$-等变）在 2007 年前后已由 Gaitsgory--Frenkel 证明：临界水平 $\hat{\mathfrak{g}}$-模范畴在球面情形下等价于非分歧 opers 上的拟凝聚层范畴。这条定理今天仍是局部理论最可靠的部分。

一般情形仍是猜想。2024 年的整体证明没有解决它，因为整体证明中的“局部系统”是整体曲线上的对象，而局部对应要求的是去心圆盘 $D^\times$ 上的 opers 与 $G((t))$ 的范畴表示之间的对应。二者的关系是“整体—局部”的纤维化，不是同一个定理。

温和分歧情形（Iwahori-等变）在本节 §10.4 被陈述为与 Springer 纤维上的拟凝聚层范畴的等价；$\mathfrak{sl}_2$ 情形已验证，一般情形仍是猜想。2007 年之后这一方向没有出现完整的证明，只有零散的特例与工具性进展。

\section*{3. q-特征标、q-Langlands 与量子几何 Langlands}

\subsection*{3.1 一个术语需要纠正}

任务书中提到的“Frenkel--Gaitsgory 的 q-特征层”，在文献中并不存在。核查结果如下：
\begin{itemize}
\item q-特征标（q-characters）的原始文献是 Frenkel--Reshetikhin，\emph{The q-characters of representations of quantum affine algebras and deformations of W-algebras}，Contemp. Math. 248 (1999) 163--205，\texttt{arXiv:math/9810055}。它是量子仿射代数表示论的工具，与“本征层”（Hecke eigensheaf）不是一回事。
\item Frenkel 与 Gaitsgory 的真实合作是\emph{局部}几何 Langlands（临界水平、仿射 Kac--Moody 代数），即 \texttt{arXiv:math/0508382}，与 q-特征标无关。
\end{itemize}
因此本章不采用“q-特征层”这一说法，而分别介绍 q-特征标与 q-Langlands 两条真实存在的线索。

\subsection*{3.2 q-特征标的 Langlands 对偶}

Frenkel--Hernandez 把 q-特征标与 Langlands 对偶联系起来（\emph{Langlands duality for finite-dimensional representations of quantum affine algebras}，Lett. Math. Phys. 96 (2011) 217--261，\texttt{arXiv:0902.0447}）。对 Kirillov--Reshetikhin 模，该对偶已证；一般情形仍是猜想。这是“部分结果”，不能写成定理。

\subsection*{3.3 真正的 q-Langlands}

文献中“q-Langlands”指的是量子几何 Langlands 的 q-形变。Aganagic--Frenkel--Okounkov 在 \emph{Quantum q-Langlands Correspondence}（Trans. Moscow Math. Soc. 79 (2018) 1--83，\texttt{arXiv:1701.03146}）中对所有 simply-laced 李代数证明了：量子仿射代数的 q-conformal blocks 等同于两个 Langlands 对偶李代数相应的形变 W-代数的块。证明用到 Nakajima quiver variety 的量子 K-理论；物理来源是六维 little string theory。

这是一条已发表、经同行评议的定理，但适用范围有限：simply-laced、q-形变、整体曲线。完整的量子几何 Langlands 仍是猜想，其系统的猜想形式见 Gaitsgory，\emph{Quantum Langlands correspondence}，\texttt{arXiv:1601.05279}（2016）。猜想的核心是水平 $\kappa$ 与 $-1/(m\kappa)$ 的互换（$m$ 为 lacing number）。表示论侧的部分结果见 Arakawa--Frenkel，\emph{Quantum Langlands duality of representations of W-algebras}，Compositio Math. 155 (2019) 2235--2262。

\section*{4. 关于“opuscule 表示”的核查}

任务书还提到“opuscule 表示与量子几何 Langlands”。本轮对 arXiv 全库与 OpenAlex 做了检索：标题或全文中出现“opuscule”的数学文献为零；命中都来自天文学或人文古典学（如 Pascal 的《Opuscules》）。也就是说，\textbf{数学文献中不存在名为 opuscule 的标准概念}，无法把它对应到任何可核对的作者或论文。

最可能的两种情况是：
\begin{enumerate}
\item 本意是 \textbf{opers}。本书第 4、8、10 章的核心对象；2024 年证明的第 IV 篇关键一步正是“几何 Langlands 等价于 generic oper structures 空间可缩”。
\item 本意是 \textbf{minuscule}（极小表示）。这是标准术语，但与量子几何 Langlands 没有直接关系。
\end{enumerate}
本章不采用“opuscule 表示”这一说法，以免向读者引入一个不存在的术语。若读者手头有使用该词的文献，欢迎指出，我们会在下一版补正。

\section*{5. W-代数、class $\mathcal S$ 与四维规范理论}

本书第 8 章证明临界水平中心同构于经典 W-代数 $\mathbf W({}^L\mathfrak{g})$，这是“Langlands 对偶群从 W-代数对偶中出现”的代数机制。2007 年之后，这条线最重要的出口是 Arakawa 的工作。

Arakawa 在预印本 \emph{Chiral algebras of class $\mathcal S$ and Moore--Tachikawa symplectic varieties}（\texttt{arXiv:1811.01577}，2018；截至 2026 年仍无期刊版）中证明：
\begin{enumerate}
\item 对任意复半单 $G$（不必 simply-laced），genus $0$ 的 class $\mathcal S$ 手征代数存在、唯一、单、共形；
\item 其 associated variety 恰为 Braverman--Finkelberg--Nakajima 用 Langlands 对偶群的仿射 Grassmannian 构造的 Moore--Tachikawa 辛簇。
\end{enumerate}
注意状态：结论以定理形式陈述，但尚未经同行评议，书稿应写“Arakawa 在预印本中证明”。

物理侧的三个来源是 AGT 对应（Alday--Gaiotto--Tachikawa 2010）、Beem 等的 4d/2d 手征代数（2015）、Moore--Tachikawa 的 2d TQFT（2012）。这些是物理提案或物理层面的已发表结果，数学上一般为猜想，不能与 Arakawa 的数学定理混为一谈。

\section*{6. Kapustin--Witten 之后}

本书引言提到 Kapustin--Witten 的物理解释。其原始文献是 \emph{Electric-magnetic duality and the geometric Langlands program}，Commun. Number Theory Phys. 1 (2007) 1--236，\texttt{arXiv:hep-th/0604151}。核心图像是：四维 $\mathcal N=4$ 超对称 Yang--Mills 的 S-对偶（$G\leftrightarrow{}^L G$）在紧 Riemann 面上拓扑扭变后，'t Hooft 线算子对应 Hecke 算子，B-膜对应本征层与局部系统。

四条主要后续：
\begin{enumerate}
\item Gukov--Witten 把分歧与表面算子纳入物理图像（\texttt{arXiv:hep-th/0612073}，2006）；
\item Witten 从六维 $(2,0)$ 理论解释几何 Langlands（\texttt{arXiv:0905.2720}，2009）；
\item Gaiotto--Witten 的解析版本，把 Hitchin Hamiltonians 与 Hecke 算子实现为 Hilbert 空间上的具体算子（\texttt{arXiv:2107.01732}；Ann. Henri Poincar\'e 25 (2024) 557--671）；
\item Frenkel--Gaiotto 的量子版本，把量子 Langlands 对偶表述为边界条件、$D$-模与共形块之间的对偶（Commun. Number Theory Phys. 14 (2020) 199--313）。
\end{enumerate}

需要强调：2024 年的数学证明与 Kapustin--Witten 的物理推导是两条独立路线。前者不依赖后者，后者也不是证明；二者在结论上会合，这正是这个领域最引人注目的地方之一。

\section*{7. 其他邻近进展}

\begin{itemize}
\item \textbf{正特征。} Gaitsgory--Raskin，\emph{Geometric Langlands in positive characteristic from characteristic zero}，\texttt{arXiv:2508.02237}（2025，预印本）。
\item \textbf{整体对应的现代形式。} Arinkin--Gaitsgory，Selecta Math. 21 (2015) 1--199，将几何 Langlands 重述为 $D\text{-}\mathrm{mod}(\mathrm{Bun}_G)\simeq\mathrm{IndCoh}_{\mathcal N}(\mathrm{LS}_{{}^L G})$；这是 2024 年五篇的技术底座。
\item \textbf{临界水平中心。} Feigin--Frenkel 1992 年的原始结果（Int. J. Mod. Phys. A 7, Suppl. 1A (1992) 197--215）仍是本书第 8 章主定理的源头，2007 年后没有被推翻，只有新的证明与推广。
\item \textbf{$p$-进侧。} Fargues--Scholz 的局部 Langlands 几何化（2018 年预印本，2021 年起陆续发表）是另一条线，与本书的圈群设定不同，但对“把 Langlands 对应几何化”这一总体纲领有重要影响。本书没有涉及，读者若对“几何化”本身的现状感兴趣，可以由此进入。
\end{itemize}

\section*{8. 本章的局限}

\begin{enumerate}
\item 五篇预印本的定理编号与逐条假设，本章只据摘要、目录与 IHES 概述，未逐页核对正文；引用具体定理时应回到原文。
\item 第 V 篇 §0.2.7 列有“Deficiencies”一节，作者自陈尚有未尽事项；本章未逐条展开。
\item Arakawa 的预印本尚无期刊版，状态可能在本版付印后变化。
\item “opuscule”的核查只覆盖 arXiv 与 OpenAlex 的标题/全文检索，不能排除某个未索引的讲义或会议论文使用过该词；如有，请读者指出。
\item 本章不涉及 Kapustin--Witten 方程本身的 PDE 分析文献。
\end{enumerate}

\section*{9. 本章引用文献}

\begin{enumerate}
\item D. Gaitsgory, S. Raskin, \emph{Proof of the geometric Langlands conjecture I: construction of the functor}, \texttt{arXiv:2405.03599}.
\item D. Arinkin et al., \emph{Proof of the geometric Langlands conjecture II: Kac--Moody localization and the FLE}, \texttt{arXiv:2405.03648}.
\item J. Campbell et al., \emph{Proof of the geometric Langlands conjecture III: compatibility with parabolic induction}, \texttt{arXiv:2409.07051}.
\item D. Arinkin et al., \emph{Proof of the geometric Langlands conjecture IV: ambidexterity}, \texttt{arXiv:2409.08670}.
\item D. Gaitsgory, S. Raskin, \emph{Proof of the geometric Langlands conjecture V: the multiplicity one theorem}, \texttt{arXiv:2409.09856}.
\item D. Gaitsgory, S. Raskin, \emph{Geometric Langlands in positive characteristic from characteristic zero}, \texttt{arXiv:2508.02237}.
\item D. Arinkin, D. Gaitsgory, \emph{Singular support of coherent sheaves and the geometric Langlands conjecture}, Selecta Math. (N.S.) 21 (2015) 1--199.
\item E. Frenkel, N. Reshetikhin, \emph{The q-characters of representations of quantum affine algebras and deformations of W-algebras}, Contemp. Math. 248 (1999) 163--205.
\item E. Frenkel, D. Hernandez, \emph{Langlands duality for finite-dimensional representations of quantum affine algebras}, Lett. Math. Phys. 96 (2011) 217--261.
\item M. Aganagic, E. Frenkel, A. Okounkov, \emph{Quantum q-Langlands correspondence}, Trans. Moscow Math. Soc. 79 (2018) 1--83.
\item D. Gaitsgory, \emph{Quantum Langlands correspondence}, \texttt{arXiv:1601.05279}.
\item T. Arakawa, E. Frenkel, \emph{Quantum Langlands duality of representations of W-algebras}, Compositio Math. 155 (2019) 2235--2262.
\item T. Arakawa, \emph{Chiral algebras of class $\mathcal S$ and Moore--Tachikawa symplectic varieties}, \texttt{arXiv:1811.01577}.
\item L. F. Alday, D. Gaiotto, Y. Tachikawa, \emph{Liouville correlation functions from four-dimensional gauge theories}, Lett. Math. Phys. 91 (2010) 167--197.
\item C. Beem et al., \emph{Infinite chiral symmetry in four dimensions}, Comm. Math. Phys. 336 (2015) 1359--1433.
\item G. Moore, Y. Tachikawa, \emph{On 2d TQFTs whose values are holomorphic symplectic varieties}, Proc. Sympos. Pure Math. 85 (2012) 191--297.
\item A. Kapustin, E. Witten, \emph{Electric-magnetic duality and the geometric Langlands program}, Commun. Number Theory Phys. 1 (2007) 1--236.
\item S. Gukov, E. Witten, \emph{Gauge theory, ramification, and the geometric Langlands program}, Current Developments in Mathematics 2006, 35--180.
\item E. Witten, \emph{Geometric Langlands from six dimensions}, CRM Proc. Lecture Notes 50, 281--310.
\item D. Gaiotto, E. Witten, \emph{Gauge theory and the analytic form of the geometric Langlands program}, Ann. Henri Poincar\'e 25 (2024) 557--671.
\item E. Frenkel, D. Gaiotto, \emph{Quantum Langlands dualities of boundary conditions, D-modules, and conformal blocks}, Commun. Number Theory Phys. 14 (2020) 199--313.
\item E. Frenkel, D. Gaitsgory, \emph{Local geometric Langlands correspondence and affine Kac--Moody algebras}, \texttt{arXiv:math/0508382}; in \emph{Algebraic Geometry and Number Theory}, Birkh\"auser, 69--260.
\item B. Feigin, E. Frenkel, \emph{Affine Kac--Moody algebras at the critical level and Gelfand--Dikii algebras}, Int. J. Mod. Phys. A 7, Suppl. 1A (1992) 197--215.
\end{enumerate}

\noindent\textbf{来源等级说明。} 本章所有 arXiv 编号与期刊卷页均经 arXiv 摘要页与 Crossref/OpenAlex 核对（访问日期 2026-11-23）。凡标“预印本”者，均未见于期刊；凡标“未核实”者，本章不引用其具体断言。没有使用知乎、百度百科、CSDN、简书、公众号、博客或视频课笔记作为依据。
